Shapley-value explanations of tabular models are usually produced by a single run of a randomised, budget-limited estimator and read as if they were deterministic. We ask when one run suffices, judged by what a reader uses: the set of top-$k$ attributes. We study Kernel SHAP and permutation SHAP (\texttt{shap} 0.51.0) on \numnCells{} (dataset, model) cells (five public datasets by four models, $d=8$ to $20$), with 50 instances and 30 runs per instance each, against exact references (enumeration) where $d\le16$ and a high-budget reference otherwise. Error falls at about the Monte-Carlo rate, but top-3 sets need more budget than the error suggests: at the smallest budgets of our grid (20 to 60 coalitions) one Kernel SHAP run recovers the reference top-3 set in \nummatchSmallKLo{}--\nummatchSmallKHi{} of runs. At the library-default budget ($2d+2048$ coalitions), which is exact for $d\le11$ and a sampling estimate beyond, one run recovers it in \numdfltMatchThreeLo{}--\numdfltMatchThreeHi{} of runs on the \numnDefaultCells{} $d>10$ cells (German Credit, ACS Employment), and the top-3 match probability is at least $0.9$ for \numdfltAdeqLo{}--\numdfltAdeqHi{} of instances; this is evidence from two datasets, one trained model per cell. On those cells two models retrained on 80\% subsamples agree on the top-3 set less often than two runs (\numstdretrainSampledLo{}--\numstdretrainSampledHi{}), under our retraining protocol. At matched compute, averaging four Kernel SHAP runs had a lower top-3 match probability than one run four times larger (pooled difference \numhmFourK{}, 95\% CI \numhmFourLoK{} to \numhmFourHiK{}, \numhmFourPairsK{} matched budget pairs); for permutation SHAP we found no difference (\numhmFourP{}, \numhmFourLoP{} to \numhmFourHiP{}). Subset checks on the library's default \texttt{l1\_reg} and on the explained output (probability or log-odds) show that both change what is reproduced.
